\chapter{Results and Discussion}\label{ch:results}

Unless stated otherwise, numbers are for the 88 primary test records and medium layer 16; ``top-1'' means that the edit's target becomes the model's first choice. Every pass bar was fixed before its run (Appendix~A, Table~\ref{tab:gates}).

\section{The Transient Edit in One Model}\label{sec:single}

The first part of the project compared three editors on 300 held-out CounterFact prompts (Figure~\ref{fig:rank1band}, Table~\ref{tab:single}). Per-feature gradient descent reached rank 1 on 271 prompts, against 161 for the Hybrid sweep given the same allowed range of 0 to 3 times per feature and 94 for the original sweep settings. It solved 110 prompts the wide sweep missed and lost none, had a lower side effect on the same prompt on 160 of the 161 prompts both solved, and took about 4 to 6.5 times less compute; part of the side-effect advantage is by construction, since the loss contains a side-effect term, and the median target probability at rank 1 is only 7.7\%. Earlier sweep studies showed that the supply of usable features, not the depth of the intervention, is the limit: drawing candidates from all prompt positions brought the target to rank 1 on 10 of 13 comparable prompts against 4 of 13 for the last token, and removing the per-feature safety filter lost 10 of 29 successes.

\begin{figure}[t]
\centering
\safegraphic[width=0.88\linewidth]{figures/gd_rank1_by_band}
\caption{Prompts brought to rank 1 (labels give the count) by starting rank of the true answer, GPT-2 Small layer 8, 300 held-out CounterFact prompts. Gradient descent is the green bars.}\label{fig:rank1band}
\end{figure}

\begin{table}[t]
\centering\small
\caption{Single-model editing on 300 held-out CounterFact prompts.}\label{tab:single}
\begin{tabular}{|p{0.50\linewidth}|p{0.22\linewidth}|p{0.15\linewidth}|}
\hline
\textbf{Editor (allowed range per feature)} & \textbf{Reached rank 1} & \textbf{Share}\\ \hline
Sweep, original settings (0.4 to 1.5 times) & 94 of 300 & 31\%\\ \hline
Sweep, wide range (0 to 3 times) & 161 of 300 & 54\%\\ \hline
Gradient descent (0 to 3 times) & \textbf{271 of 300} & \textbf{90\%}\\ \hline
\end{tabular}
\end{table}

Reaching rank 1 proved a weak test (Figure~\ref{fig:additive}). An additive edit that switches on silent features pushed random unrelated words to rank 1 on 20 of 20 prompts (the multiplier edit: 1 of 20), at a median size of 126 to 192\% of the residual norm against 39\% for multipliers; on 20 unrelated prompts it changed the top word 24\% of the time (multipliers: 6\%) and its transfer to reworded prompts was as large for a random word as for the true answer. Reachability by steering is therefore not evidence about what a model knows, and an early classification of some facts as ``absent from the model'' was withdrawn. This is why every later result is judged against controls.

\begin{figure}[t]
\centering
\safegraphic[width=0.80\linewidth]{figures/fig_additive}
\caption{An additive edit forces random words to rank 1 (a) at sizes comparable to or larger than the residual stream itself (b); the multiplier edit does not.}\label{fig:additive}
\end{figure}

\section{The Channel: Translator Quality and a Positive Control}\label{sec:channel}

A ridge map fitted on wikitext-2 keeps the receiver working (Table~\ref{tab:translator}): stitching one model's first layers to the other's last layers through the map recovers 91\% (small to medium) and 98\% (medium to small) of the receiver's own quality; in the second case the stitched loss is 4.31 against 4.24 for small alone and 3.96 for medium. The map copies which word changed (cosine of the translated and true difference at the changed word 0.63 to 0.76 over eight layer pairs) but the summary at the last position only weakly (0.24 to 0.31); only the two medium to small pairs reached the 0.3 line fixed for the translator gate, so the translator passes that gate for Import and not for Export.

\begin{table}[t]
\centering\small
\caption{Ridge translators for the two directions studied (wikitext-2, held-out tokens).}\label{tab:translator}
\begin{tabular}{|p{0.30\linewidth}|p{0.15\linewidth}|p{0.17\linewidth}|p{0.22\linewidth}|}
\hline
\textbf{Direction} & \textbf{$R^2$ per dimension} & \textbf{Stitching recovers} & \textbf{Difference cosine at the last position}\\ \hline
small layer 8 $\to$ medium layer 16 (Export) & 0.58 & 91\% & 0.28\\ \hline
medium layer 16 $\to$ small layer 8 (Import) & 0.66 & 98\% & 0.31\\ \hline
\end{tabular}
\end{table}

As a positive control we swap one country for another in a template sentence. The change is computed as a difference of the source model's states, translated, and added to the receiver, and counted as a hit when the receiver's top answer becomes the swapped capital. The translated change succeeds on 55\% (Export) and 33\% (Import) of test cases, against at most 0.4\% for a random vector, a random map and another pair's difference (Figure~\ref{fig:swap}). The control is mostly a channel check at the changed word's position; its test cases share countries and are not independent, and it is neither an edit nor a fact.

\begin{figure}[t]
\centering
\safegraphic[width=0.92\linewidth]{figures/r_swap}
\caption{Country-swap control: the translated difference moves the receiver far above the best control (tick).}\label{fig:swap}
\end{figure}

\section{Export with the Text-Fitted Map}\label{sec:export}

For each of the 88 primary test records the edit tuned in small (200 multipliers, rank 1 in small) is relayed to medium through the ridge map, scaled by a dose chosen on dev, and added at layer 16 (Figure~\ref{fig:doses}, Table~\ref{tab:export}). The effect is real, graded and specific. The target's median rank in medium falls from 142 to 14 (23 at layer 12), it becomes the top answer on 8\% of records (5\% at layer 12), reworded prompts favour the new answer on 52\% against 36\% unchanged, and neighbours lose only 4 points. A random vector and another record's edit do nothing (0\% top-1, ranks 157 and 128). A plain word push along the target's output direction makes the target the top answer on 83\% of records, but only 23\% of neighbours keep the true answer, and stating the fact in medium's prompt lifts rewordings to 97\% at the cost of 23 points of neighbours (Figure~\ref{fig:side}). The pass bar fixed beforehand (top-1 of at least 15\% and 10 points above the strongest control) was not met. The edit moves medium toward the target without damaging its other answers, but it rarely does so cleanly enough to reach rank 1.

\begin{figure}[t]
\centering
\safegraphic[width=\linewidth]{figures/r_export_doses}
\caption{Export of the edit tuned in small, small layer 8 to medium layer 16, 88 test records: rank (a) and top-1 rate (b) against the dose. The word push uses the answer and is the comparison, not a method.}\label{fig:doses}
\end{figure}

\begin{table}[t]
\centering\small
\caption{Export of the edit through the ridge map (88 primary test records; dose chosen on dev). Unchanged medium: median rank 142.}\label{tab:export}
\begin{tabular}{|p{0.38\linewidth}|p{0.10\linewidth}|p{0.12\linewidth}|p{0.12\linewidth}|p{0.12\linewidth}|}
\hline
\textbf{Arm} & \textbf{Top-1} & \textbf{Median rank} & \textbf{Rewordings} & \textbf{Neighbours}\\ \hline
Medium unchanged & 0\% & 142 & 36\% & 67\%\\ \hline
\textbf{Translated edit} & \textbf{8\%} & \textbf{14} & \textbf{52\%} & \textbf{63\%}\\ \hline
Random vector & 0\% & 157 & 38\% & 66\%\\ \hline
Another record's edit & 0\% & 128 & 36\% & 65\%\\ \hline
Edit toward a random word (own word) & 3\% & 31 & -- & --\\ \hline
Word push (uses the answer) & 83\% & 1 & 82\% & 23\%\\ \hline
Fact stated in medium's prompt & -- & -- & 97\% & 44\%\\ \hline
Small with the edit itself & -- & -- & 67\% & 54\%\\ \hline
\end{tabular}
\end{table}

\begin{figure}[t]
\centering
\safegraphic[width=\linewidth]{figures/r_export_side}
\caption{What the export does to reworded prompts and to neighbouring facts. Grey bars are references. The word push and the fact in the prompt damage neighbours far more than the translated edit.}\label{fig:side}
\end{figure}

\section{Searching for a Better Translator}\label{sec:search}

Because the effect is real but weak, we tried seven changes, each with its pass rule written first (Table~\ref{tab:ladder}, Figure~\ref{fig:ladder}).

\textbf{Layer and capacity.} Injecting at layers 4, 8, 12, 16 and 20 gave 0, 5, 5, 8 and 1\%; splitting the dose over several layers did no better (best set 7\%), so the layer is not the limit. A neural map fitted on the same text fits it much better ($R^2$ +0.09 to +0.11) but exports no better (8\%): a better fit to ordinary text does not help.

\textbf{A map trained on edits.} Training the linear map so that medium's answers move as small's do (method 3 of Chapter~\ref{ch:method}) on 289 successful edits roughly doubled the result, from 8\% to 18\% top-1 and from median rank 14 to 7 (paired $p=5\times10^{-9}$), still specific to the record (another record's edit: 0\%). It did not reach the 50\% bar. More training edits (25, 50, 100\%: 13, 13, 20\%), longer training and a direct rank-1 term in the loss (16 to 21\% for every weight) left it at about 20\%. One step produced a misleading figure: with 40 epochs the dose chosen on 28 dev records by top-1 was 1.0 and gave 2\% against 20\%; a post-hoc check showed that at dose 2.0 both lengths give 20\%, and the dose is chosen by the dev log-rank gain from then on.

\textbf{Why it plateaus.} With the rank-1 term the map learned its 289 training targets (95\% top-1 on them) but not unseen words (16\%): it memorises. We then gave it 1,221 more edits toward 759 additional random target words (the 289 real edits cover only 116). Memorisation vanished (training minus test top-1: 12 points down to 3) but unseen records got \emph{worse}, 20\% to 14\% (Figure~\ref{fig:gap}); training top-1 fell as well while the training loss stayed the same, because the loss asks medium to copy small's shift, not to reach rank 1. Variety of target words was therefore not the limit.

\textbf{Per-sentence tuning is a ceiling.} Re-tuning the 200 multipliers against medium's output for each sentence reaches the top answer on 95\% of records (84\% on all 150), with the same neighbour cost as the source edit inside small (8 points against 9). The channel can carry a rank-1 edit; what fails is the compatibility of an edit tuned for small with medium. The solution found is only partly in the direction of the original edit (cosine 0.51 at the last position against 0.02 by chance, and 0.71 of its size), so no map that merely transforms the given edit can supply what is missing. This ceiling is not a transfer, since medium's output is used for every sentence, and it also reaches 83\% for a random word.

\begin{table}[t]
\centering\small
\caption{The search for a better translator (88 primary test records, same edits, controls and judging throughout).}\label{tab:ladder}
\begin{tabular}{|p{0.07\linewidth}|p{0.38\linewidth}|p{0.11\linewidth}|p{0.11\linewidth}|p{0.17\linewidth}|}
\hline
\textbf{Step} & \textbf{What changed} & \textbf{Top-1} & \textbf{Rank} & \textbf{Pass bar}\\ \hline
D3 & Ridge map fitted on text (baseline) & 8\% & 14 & fail\\ \hline
-- & Neural map fitted on text & 8\% & 14 & no better\\ \hline
D5 & Several injection layers & 7\% & 11 & fail\\ \hline
D7 & Map trained on 289 edits & 18\% & 7 & fail\\ \hline
D8--D10 & More data, longer training, rank-1 loss term & 13--21\% & 7 & fail\\ \hline
D13 & + 1,221 edits toward 759 new words & 14\% & 9 & fail\\ \hline
D6 & Dials re-tuned per sentence (ceiling) & 95\% & 1 & neighbours fail\\ \hline
\end{tabular}
\end{table}

\begin{figure}[t]
\centering
\safegraphic[width=\linewidth]{figures/r_ladder}
\caption{Top-1 of every export variant on the same 88 test records. Blue: map fitted on text; green: map trained on edits; red: per-sentence tuning against medium's output (a ceiling, not a transfer).}\label{fig:ladder}
\end{figure}

\begin{figure}[t]
\centering
\safegraphic[width=0.92\linewidth]{figures/r_training_gap}
\caption{More and more varied training edits: top-1 on the map's own training prompts and on unseen test records (3 seeds, dots). The gap closes, but the test result does not rise.}\label{fig:gap}
\end{figure}

\section{What Is Carried: a Fact or a Push?}\label{sec:push}

For the edit inside small, a pre-stated rule gave ``read-out steering'': the gain on reworded prompts was 1.04 against 0.82 for a rank-matched random-word edit ($p=0.032$), and an edit restricted to the subject's tokens succeeded on 12\% of records against 59\% for all positions; the evidence is weak either way. For Export, each arm's rewordings gain must be measured against its \emph{own} unchanged level (a random word starts lower, at 25\% against 36\%). Doing so (Figure~\ref{fig:push}), D6's dials gain 24.4 points for the real target and 17.0 for a random word (difference +7.4, 95\% interval 0.0 to +15.3, $p=0.031$), and the map trained on edits gains 19.9 and 18.2 ($p=0.36$); the random word also loses more neighbours (12.3 against 7.3 points). Neither difference passes the $p<0.01$ line fixed beforehand. An earlier statement that the real target's gain was four times the random word's compared the wrong baselines and was withdrawn. What carries to reworded prompts is therefore a push toward the targeted word, not something specific to the fact.

\begin{figure}[t]
\centering
\safegraphic[width=0.92\linewidth]{figures/r_push}
\caption{Gain on reworded prompts, each word against its own unchanged level: the real target and a random word gain about the same.}\label{fig:push}
\end{figure}

\section{Does Small's SAE Matter for Steering Medium?}\label{sec:sae}

We repeated the per-sentence tuning with 200 random directions in place of the 200 SAE decoder directions (rows scaled to the SAE norms, the same firing pattern, loss and settings). The SAE directions reached 95\% top-1 and the random ones 81\% (62\% against 83\% for a random-word target), a difference of 15 points that meets the pre-set rule that the SAE basis ``matters''; rewordings (60\% against 58\%) and neighbours (59\% against 57\%) are almost the same (Figure~\ref{fig:sae}). The SAE makes the tuning more effective but is not what makes steering possible. The random basis still uses the SAE's activation pattern, so ``no SAE at all'' is untested, and one random draw was used.

\begin{figure}[t]
\centering
\safegraphic[width=0.92\linewidth]{figures/r_sae}
\caption{The tuning of Figure~\ref{fig:ladder}, step D6, with small's SAE directions and with 200 random directions (88 test records).}\label{fig:sae}
\end{figure}

\section{Import: Medium's State into Small}\label{sec:import}

Of the 21,789 facts, medium alone gets 9.9\% right and small 8.3\%; 780 facts are right in medium and wrong in small (for 71\% of them the true answer is small's second to fifth choice). \textbf{Full replacement} of small's layer-8 state by medium's translated state rescues 262 of these 780 (33.6\%; controls: wrong-sentence state 0.8\%, average state 0\%, random map 0\%), but breaks 22\% of the facts both models knew and 70\% of those only small knew, so small's accuracy falls from 8.30\% to 7.17\%; a filter through small's SAE rescued only 6.3\%. The pre-set bar (50\% rescue, accuracy at least half-way to medium, 90\% of known facts kept) was not met, although the controls show that real information crosses.

A \textbf{softer mix} was then tried, with its rule written before the build: $\gamma\in\{0.1,\dots,1\}$ and gating by medium's confidence, the setting chosen on one half of the facts and reported on the other (half B, 10,895 facts). Mixing in half of medium's state ($\gamma=0.5$) raises small's accuracy from 8.5\% to 9.5\% (medium alone: 10.2\%), gaining 162 facts and losing 53 ($p=2.5\times10^{-14}$), keeps 98\% of the facts both knew and 80\% of those only small knew, and rescues 28.6\% of the facts small gets wrong (Table~\ref{tab:import}, Figure~\ref{fig:import}); gating by confidence did not beat the plain mix. This meets the rule (gain of at least 0.5 points, 90\% of known facts kept, $p<0.01$) and is the first net-positive Import result. It is a small gain (one point), it needs medium to run on the same sentence, the old rescue bar of 50\% is still not met, and the controls of the replacement test were not rerun for the mix.

\begin{table}[t]
\centering\small
\caption{Import on the held-out half B of the facts (10,895 facts): replacement and the 50/50 mix against the two models alone.}\label{tab:import}
\begin{tabular}{|p{0.37\linewidth}|p{0.10\linewidth}|p{0.11\linewidth}|p{0.14\linewidth}|p{0.11\linewidth}|}
\hline
\textbf{Measure} & \textbf{Small} & \textbf{Medium} & \textbf{Replace} ($\gamma=1$) & \textbf{Mix} ($\gamma=0.5$)\\ \hline
Accuracy over all facts & 8.5\% & 10.2\% & 7.4\% & \textbf{9.5\%}\\ \hline
Facts small got wrong, now right & -- & -- & 31.9\% & 28.6\%\\ \hline
Facts both knew, still right & 100\% & -- & 78.9\% & 98.3\%\\ \hline
Facts only small knew, still right & 100\% & -- & 33.3\% & 80.5\%\\ \hline
Agreement with medium's top answer & 71.0\% & 100\% & 74.5\% & 76.2\%\\ \hline
\end{tabular}
\end{table}

\begin{figure}[t]
\centering
\safegraphic[width=\linewidth]{figures/r_import}
\caption{Import by mixing small's own state with medium's translated state. (a) What the mix keeps and fixes. (b) Net accuracy: best at a 50/50 mix, below small alone under full replacement.}\label{fig:import}
\end{figure}

\section{Discussion and Threats to Validity}\label{sec:discussion}

\textbf{What the evidence supports.} A linear channel between GPT-2 Small and Medium exists and keeps the receiver working. The change an edit makes in small moves medium's answers specifically, in a graded way, and without damaging its neighbours, but every translator that does not look at medium's output stops at about 20\% top-1, far below the 50\% bar; looking at medium for every sentence reaches 95\%. Together with the controls this puts the limit on the \emph{edit's compatibility}, not on the channel, the layer, the translator's capacity or the amount of data. What crosses is a push toward the targeted word and not a fact. In the other direction a half-and-half mix gives a small but statistically clear gain.

\textbf{Threats to validity.} (1) One model pair, one dataset (single-token CounterFact facts) and one SAE; nothing here says the results transfer to other models, layers or relations. (2) 88 primary test records, so one record is about one point, and training seeds share their data, so ranges understate the noise. (3) Choices of layer, dose and epoch rest on 28 to 50 dev records; the dose rule was changed once after a failure and the change is labelled post-hoc. (4) Several secondary results are descriptive diagnostics with their own rules, not method gates (the soft Import, the SAE control). (5) The additive edit, generated text, other model pairs and comparison with ROME and MEMIT were not tested. (6) The map follows earlier model-stitching work \cite{r44} and is not claimed as new.
